\chapter{Introduction}\label{ch:introduction}

\section{Background and Motivation}\label{sec:background}

Laptop selection is a recurrent decision for both individuals and organisations because the
device combines a relatively short replacement cycle with a rapidly changing catalogue of
brands, models, and configurations. Historical shipment evidence illustrates the scale of this
decision, with \textcite{sonmez_cakir_determination_2020} reporting 161.6 million laptop units
sold in 2017 and \textcite{gaur_application_2023} documenting the further expansion associated
with remote work and study in 2020--2021. The relevance of these figures lies less in the
temporary pandemic surge than in the persistent replacement problem created by a product life
commonly estimated at approximately five years \parencite{sonmez_cakir_determination_2020}.

Replacement does not reduce the decision to choosing the newest or fastest device. A prospective
buyer must compare attributes measured on different scales, including price, processor
performance, memory, storage, weight, battery capacity, and brand-related considerations. These
attributes also pull against one another rather than varying independently. A lower price
constrains hardware capability because a manufacturer can only fit so much component cost within
a given margin, and a larger battery adds weight because the energy density of lithium-ion cells
is bounded, so additional capacity has to come from a physically heavier pack. The proliferation
of configurations therefore increases the information available to the buyer without necessarily
making the choice more coherent. An explicit decision model is useful because it records how
competing attributes are weighted and how each alternative performs against the same criteria.

A problem of this form falls within Multi-Criteria Decision-Making (MCDM). The systematic review
by \textcite{basilio_systematic_2022} documents the application of MCDM across a wide range of
domains, with the Analytic Hierarchy Process (AHP) and the Technique for Order of Preference by
Similarity to Ideal Solution (TOPSIS) among the most frequently applied methods. Their prevalence
is also visible in the bibliometric analysis of \textcite{zyoud_bibliometric-based_2017}. AHP and
TOPSIS perform different analytical roles. AHP derives a priority vector from pairwise
judgements, whereas TOPSIS ranks alternatives by their relative distance from positive and
negative ideal solutions. Used together, they separate the elicitation of preferences from the
evaluation of products.

\section{Problem Statement}\label{sec:problem}

This study examines the decision faced by an office worker in Vietnam who must select a laptop
for work. The study does not impose a universal price band, because budgets vary across individual
and organisational purchases. Price is instead treated as a cost criterion within the decision
model; in an applied procurement setting, a buyer may additionally use a budget ceiling as an
eligibility rule before ranking the remaining alternatives. This distinction avoids assuming that
one expenditure range represents every office buyer while preserving the role of affordability in
the trade-offs that follow.

The office-worker context also affects how criteria are interpreted. Prior studies
identify recurring attributes such as price, processing capability, brand, battery life, and
portability, but their priorities are derived from different products, markets, and user
populations \parencite{sonmez_cakir_determination_2020,liao_determinants_2022,weng_siew_lam_evaluation_2023}.
Those priorities cannot be assumed to transfer unchanged to Vietnamese working conditions, where
local factors such as the ready availability of mains power in offices or the prevalence of
motorbike commuting may change how much a given attribute contributes to the decision. The
criteria set must first be examined against local experience and then screened by the intended
user population before its relative weights are elicited.

Previous research demonstrates that hybrid MCDM models can rank laptops and related electronic
devices. In particular, \textcite{gaur_application_2023} combine AHP and the Best--Worst Method
for weighting with TOPSIS for laptop ranking, while \textcite{weng_siew_lam_evaluation_2023}
apply an integrated AHP--TOPSIS model to mobile phones. The unresolved issue is therefore not
whether AHP and TOPSIS can be combined. It is how a laptop-ranking model can maintain a traceable
connection between literature-derived criteria, local expert interpretation, wider user
screening, preference elicitation, and a date-stamped set of product alternatives.

\section{Research Gap and Contributions}\label{sec:gap-contrib}

The literature reveals a separation between three forms of evidence. Laptop-selection studies
identify criteria and sometimes rank alternatives, but commonly begin from a predetermined
criteria set. Consumer studies provide evidence about attributes that influence purchase, yet do
not necessarily translate those attributes into an auditable ranking model. Context-specific
qualitative evidence can reveal local constraints, but expert observations alone do not establish
that the wider target population assigns the same importance to each criterion. Few studies join
these forms of evidence in a staged process that distinguishes criterion identification,
population screening, preference weighting, and alternative ranking.

This study addresses that design gap through two contributions. First, it develops a
context-calibrated criteria set for office workers in Vietnam by combining the international
literature, interviews with five experts occupying different positions in laptop procurement, and
a screening survey of 84 respondents. This sequence makes agreements and divergences between
published evidence, professional judgement, and user ratings visible. Second, it implements a
reproducible AHP--TOPSIS application in which one representative expert supplies the pairwise AHP
judgements and TOPSIS ranks twelve internationally recommended laptop models using price and
configuration data recorded from Amazon.com on 12 July 2026. The resulting ranking is an
evaluation of an international candidate snapshot for a Vietnamese office-worker decision
context; it is not presented as a census of the Vietnamese retail market.

\section{Objectives and Research Questions}\label{sec:objectives}

The study has three objectives. The first is to establish criteria appropriate for evaluating
office laptops in the Vietnamese context. The second is to estimate the relative importance of
the retained criteria through AHP. The third is to apply those weights in TOPSIS to rank a defined
set of laptop alternatives. These objectives are expressed through the following research
questions:

\begin{enumerate}
  \item \textbf{RQ1:} Which criteria are appropriate for evaluating laptops for office workers in Vietnam?
  \item \textbf{RQ2:} What weights do these criteria carry when determined through AHP?
  \item \textbf{RQ3:} Under a simulated purchase scenario, which laptop model is the most suitable choice under the established criteria and weights when ranked with TOPSIS?
\end{enumerate}

\section{Scope and Structure of the Report}\label{sec:scope}

The unit of analysis is a laptop considered for office work by a Vietnamese buyer. The study
excludes specialised gaming and graphics-workstation requirements and does not claim to represent
all user segments or retail channels. The twelve alternatives originate from current
recommendations published by seven professional review outlets; their price and configuration
data are taken from Amazon.com because most of the newly released international models were not
available in the Vietnamese catalogues reviewed at the collection date. Market data are therefore
treated as a dated snapshot rather than a stable property of the alternatives. The study further
uses a single representative expert for AHP weighting, so the resulting priority vector reflects
one coherent expert judgement rather than a group-AHP consensus.

The remainder of the report is organised into five chapters. \Cref{ch:literature} establishes the
theoretical foundations of MCDM, AHP, and TOPSIS and synthesises prior laptop and electronics
selection studies. \Cref{ch:methodology} explains the staged criteria-development process,
preference elicitation, alternative identification, and computational procedures.
\Cref{ch:results} reports the screening, weighting, and ranking results. \Cref{ch:discussion}
interprets those findings against prior evidence, and \Cref{ch:conclusion} summarises the
contributions, limitations, and directions for further research.
